\documentclass[tikz,border=25pt]{standalone}
\usepackage[utf8]{inputenc}
\usepackage[russian]{babel}
\usepackage{tikz}
\usetikzlibrary{shapes.geometric, arrows.meta, positioning, calc}

\begin{document}
\begin{tikzpicture}[
    x=1cm, y=1cm,
    >=Stealth,
    font=\sffamily\scriptsize,
    % Component Styles
    regnode/.style={
        draw=black!85, rectangle, rounded corners=3pt, minimum width=1.1cm, minimum height=0.75cm, 
        line width=1.1pt, fill=white, font=\bfseries\scriptsize, align=center
    },
    block/.style={
        draw=black!85, rectangle, rounded corners=3pt, line width=1.1pt, fill=white, 
        font=\bfseries\scriptsize, align=center
    },
    ctrlbox/.style={
        draw=red!80!black, rectangle, rounded corners=7pt, line width=1.4pt,
        fill=red!5, align=center
    },
    membox/.style={
        draw=cyan!75!black, rectangle, rounded corners=6pt, line width=1.3pt,
        fill=cyan!6, align=center
    },
    % Wire Styles
    databus/.style={->, line width=1.3pt, draw=black!90},
    databus_line/.style={line width=1.3pt, draw=black!90},
    ctrlbus/.style={->, line width=1.0pt, draw=red!80!black, dashed},
    ctrlbus_line/.style={line width=1.0pt, draw=red!80!black, dashed},
    bidir/.style={<->, line width=1.2pt, draw=teal!80!black}
]

    % ========================================================
    % 1. РЕГИСТРЫ R0 - R7 (16-BIT GPR) (Y = 12.0)
    % ========================================================
    \foreach \i/\x in {0/0.6, 1/2.0, 2/3.4, 3/4.8, 4/6.2, 5/7.6, 6/9.0, 7/10.4} {
        \node[regnode] (R\i) at (\x, 12.0) {\textbf{R\i}};
        \coordinate (R\i_in) at (\x - 0.25, 12.38);
        \coordinate (R\i_we) at (\x + 0.25, 12.38);
        \coordinate (R\i_out) at (\x, 11.62);
    }

    % ========================================================
    % 2. MUX 1 (REG MUX) (Y = 9.7 .. 10.6, Center X = 5.5)
    % ========================================================
    \draw[fill=gray!12, draw=black!85, line width=1.2pt] 
        (-0.2, 10.6) -- (11.2, 10.6) -- (7.9, 9.7) -- (3.1, 9.7) -- cycle;
    \node[align=center] at (5.5, 10.15) {\textbf{\normalsize MUX 1}};

    % Спуски от каждого регистра в MUX 1
    \foreach \i in {0,...,7} {
        \draw[databus] (R\i_out) -- (R\i_out |- 0, 10.6);
    }

    % ========================================================
    % 3. ЛЕВАЯ КОЛОНКА: DEST-DMUX, ACCA, ACCB, ALU (16-BIT)
    % ========================================================
    % DEST-DMUX (Center X = 2.2, Y = 5.9 .. 6.9)
    \draw[fill=gray!12, draw=black!85, line width=1.1pt] 
        (1.3, 6.9) -- (3.1, 6.9) -- (3.6, 5.9) -- (0.8, 5.9) -- cycle;
    \node[align=center] at (2.2, 6.4) {\textbf{\scriptsize DEST-DMUX}};

    % Аккумуляторы ACCA и ACCB (Y = 4.9)
    \node[regnode, fill=blue!6, draw=blue!75!black, minimum width=1.2cm, minimum height=0.8cm] 
        (ACCA) at (1.3, 4.9) {\textbf{ACCA}};
    \node[regnode, fill=blue!6, draw=blue!75!black, minimum width=1.2cm, minimum height=0.8cm] 
        (ACCB) at (3.1, 4.9) {\textbf{ACCB}};

    % Связи от DMUX к аккумуляторам
    \draw[databus] (1.3, 5.9) -- (ACCA.north);
    \draw[databus] (3.1, 5.9) -- (ACCB.north);
    \node[font=\tiny\bfseries, fill=white, inner sep=1pt] at (1.05, 5.55) {[0]};
    \node[font=\tiny\bfseries, fill=white, inner sep=1pt] at (3.35, 5.55) {[1]};

    % ALU (Center X = 2.2, Y = 2.0 .. 3.6)
    \draw[fill=blue!8, draw=blue!85!black, line width=1.3pt]
        (0.5, 3.6) -- (1.8, 3.6) -- (2.2, 3.2) -- (2.6, 3.6) -- (3.9, 3.6) -- (3.2, 2.0) -- (1.2, 2.0) -- cycle;
    \node[align=center, text=blue!90!black] at (2.2, 2.7) {\textbf{\large ALU}};

    % Входы в ALU от аккумуляторов
    \draw[databus] (ACCA.south) -- (1.3, 3.6);
    \draw[databus] (ACCB.south) -- (3.1, 3.6);
    \node[font=\tiny\bfseries, left] at (1.3, 4.1) {A [15:0]};
    \node[font=\tiny\bfseries, right] at (3.1, 4.1) {B [15:0]};

    % ========================================================
    % 4. ЦЕНТРАЛЬНАЯ КОЛОНКА: MUX 2 (MEM MUX) И FLAG REG
    % ========================================================
    % MUX 2: ПЕРЕВЕРНУТ! (Узкий верх Y = 5.8, широкий низ Y = 4.4, Center X = 6.4)
    \draw[fill=yellow!18, draw=orange!80!black, line width=1.3pt]
        (5.6, 5.8) -- (7.2, 5.8) -- (7.8, 4.4) -- (5.0, 4.4) -- cycle;
    \node[align=center, text=orange!90!black] at (6.4, 5.25) {\textbf{\normalsize MUX 2}};

    % Бейджи входов MUX 2 ВНУТРИ трапеции по нижнему краю
    \node[font=\tiny\bfseries, fill=white, draw=orange!80!black, rounded corners=1pt, inner sep=1.2pt] at (5.6, 4.68) {[0] ALU};
    \node[font=\tiny\bfseries, fill=white, draw=orange!80!black, rounded corners=1pt, inner sep=1.2pt] at (6.4, 4.68) {[1] REG};
    \node[font=\tiny\bfseries, fill=white, draw=orange!80!black, rounded corners=1pt, inner sep=1.2pt] at (7.2, 4.68) {[2] CTRL};

    % Флаговый регистр (Center X = 6.4, Y = 1.2, W = 3.0cm, H = 0.8cm)
    \node[block, minimum width=3.0cm, minimum height=0.8cm, fill=teal!6, draw=teal!80!black] 
        (FLAGS) at (6.4, 1.2) {\textbf{FLAGS}\\[1pt]\scriptsize$[$ Z, C, N, V $]$};

    % ========================================================
    % 5. ПРАВАЯ КОЛОНКА: КОНТРОЛЛЕР И ПАМЯТЬ
    % ========================================================
    % Контроллер (X = 11.4 .. 16.8, Y = 4.2 .. 9.2, Center X = 14.1)
    \node[ctrlbox, minimum width=5.4cm, minimum height=5.0cm] 
        (CTRL) at (14.1, 6.7) {};
    
    % Заголовок контроллера
    \node[font=\bfseries\normalsize, text=red!90!black] at (14.1, 8.7) {CONTROL UNIT};

    % Порты на левой стенке Контроллера (X = 11.4)
    \node[draw=red!40!black, fill=white, rounded corners=2pt, font=\tiny\bfseries, inner sep=2pt, right] 
        at (11.5, 8.3) {IR [15:0]};
    \node[draw=red!70!black, fill=white, rounded corners=2pt, font=\tiny\bfseries, inner sep=2pt, text=red!80!black, right] 
        at (11.5, 7.3) {CTRL ALU/ACC};
    \node[draw=red!70!black, fill=white, rounded corners=2pt, font=\tiny\bfseries, inner sep=2pt, text=red!80!black, right] 
        at (11.5, 6.2) {SEL\_WB [1:0]};
    \node[draw=teal!70!black, fill=white, rounded corners=2pt, font=\tiny\bfseries, inner sep=2pt, text=teal!90!black, right] 
        at (11.5, 5.2) {FLAGS};
    \node[draw=orange!80!black, fill=white, rounded corners=2pt, font=\tiny\bfseries, inner sep=2pt, text=orange!90!black, right] 
        at (11.5, 4.5) {MEM/IMM};

    % Память системы (X = 11.4 .. 16.8, Y = 0.8 .. 3.4, Center X = 14.1)
    \node[membox, minimum width=5.4cm, minimum height=2.6cm] 
        (MEM) at (14.1, 2.1) {\textbf{\normalsize MEMORY (64 KB)}\\[3pt]\tiny ADDR [15:0] • DATA [15:0] • RD/WR};

    % Прямые вертикальные шины между Контроллером и Памятью
    \draw[databus, draw=blue!80!black] (12.6, 4.2) -- (12.6, 3.4) node[midway, fill=white, inner sep=1pt, font=\tiny\bfseries, text=blue!80!black] {ADDR};
    \draw[databus] (14.1, 4.2) -- (14.1, 3.4) node[midway, fill=white, inner sep=1pt, font=\tiny\bfseries] {WR\_DATA};
    \draw[bidir, draw=red!80!black] (15.6, 3.4) -- (15.6, 4.2) node[midway, fill=white, inner sep=1pt, font=\tiny\bfseries, text=red!80!black] {RD / WR};

    % ========================================================
    % 6. РАЗВОДКА ШИН ДАННЫХ (DATA PATH 16-BIT)
    % ========================================================
    % Спуск от MUX 1: до распределительной линии Y = 8.3
    \draw[databus_line] (5.5, 9.7) -- (5.5, 8.3);
    \filldraw (5.5, 8.3) circle (2pt);

    % Ветвь A: MUX 1 -> DEST-DMUX
    \draw[databus] (5.5, 8.3) -| (2.2, 6.9) node[pos=0.25, above, font=\tiny\bfseries] {[15:0]};

    % Ветвь C: MUX 1 -> Контроллер (порт IR)
    \draw[databus] (5.5, 8.3) -- (11.4, 8.3) node[pos=0.45, above, font=\tiny\bfseries] {REG\_DATA [15:0]};

    % Ветвь B: MUX 1 -> MUX 2 (Вход [1] REG - входит СНИЗУ)
    \filldraw (4.6, 8.3) circle (1.8pt);
    \draw[databus_line] (4.6, 8.3) -- (4.6, 3.6);
    % Горизонталь по Y = 3.6: мост над шиной ALU (X = 5.6) и вход снизу в [1] (6.4, 4.4)
    \draw[databus_line] (4.6, 3.6) -- (5.4, 3.6);
    \draw[databus_line] (5.4, 3.6) arc[start angle=180, end angle=0, radius=0.2] -- (6.4, 3.6);
    \draw[databus] (6.4, 3.6) -- (6.4, 4.4);
    \node[font=\tiny\bfseries, fill=white, inner sep=1.2pt, below] at (5.0, 3.55) {[15:0]};

    % ALU -> MUX 2 (Вход [0] ALU - входит СНИЗУ)
    \draw[databus] (3.2, 2.5) -- (5.6, 2.5) -- (5.6, 4.4);
    \node[font=\tiny\bfseries, fill=white, inner sep=1pt] at (3.9, 2.75) {ALU\_RES [15:0]};

    % CONTROLLER -> MUX 2 (Вход [2] CTRL)
    \draw[databus_line] (11.4, 4.5) -- (10.3, 4.5) -- (10.3, 3.6) -- (9.8, 3.6);
    \draw[databus_line] (9.8, 3.6) arc[start angle=0, end angle=180, radius=0.2] -- (7.2, 3.6);
    \draw[databus] (7.2, 3.6) -- (7.2, 4.4);
    \node[font=\tiny\bfseries, fill=white, inner sep=1pt] at (10.8, 4.7) {MEM/IMM [15:0]};

    % ALU -> FLAG REG
    \draw[->, line width=1.1pt, draw=teal!80!black] (2.2, 2.0) |- (FLAGS.west);
    \node[font=\tiny\bfseries, fill=white, inner sep=1pt, text=teal!90!black] at (3.5, 1.2) {FLAGS};

    % ========================================================
    % 7. FLAG REG <-> КОНТРОЛЛЕР (ДВУСТОРОННЯЯ СВЯЗЬ)
    % ========================================================
    \draw[bidir] (7.9, 1.2) -- (9.6, 1.2) -- (9.6, 5.2) -- (11.4, 5.2);
    \node[font=\tiny\bfseries, fill=white, inner sep=1pt, text=teal!90!black] at (10.4, 5.45) {FLAGS};

    % ========================================================
    % 8. ШИНА ОБРАТНОЙ ЗАПИСИ (ВЫХОД СВЕРХУ MUX 2 -> РЕГИСТРЫ)
    % ========================================================
    \draw[databus_line] (6.4, 5.8) -- (6.4, 9.0) -- (5.7, 9.0) 
                        arc[start angle=0, end angle=180, radius=0.2] 
                        -- (-0.8, 9.0) -- (-0.8, 13.5) -- (10.9, 13.5);
    \node[font=\tiny\bfseries, fill=black!85, text=white, inner sep=2.5pt, rounded corners=2pt] 
        at (2.2, 9.0) {WB\_DATA [15:0]};
    
    \foreach \i in {0,...,7} {
        \draw[databus] (R\i_in |- 0, 13.5) -- (R\i_in);
        \filldraw (R\i_in |- 0, 13.5) circle (1.8pt);
    }

    % ========================================================
    % 9. СИГНАЛЫ УПРАВЛЕНИЯ
    % ========================================================
    % 1. К MUX 1: SEL_REG [2:0]
    \draw[ctrlbus] (12.8, 9.2) |- (7.9, 10.15);
    \node[font=\tiny\bfseries, text=red!80!black, fill=white, inner sep=1pt] at (12.8, 9.7) {SEL\_REG [2:0]};

    % 2. К MUX 2: SEL_WB [1:0]
    \draw[ctrlbus] (11.4, 6.2) -| (8.5, 5.3) -- (7.4, 5.3);
    \node[font=\tiny\bfseries, text=red!80!black, fill=white, inner sep=1pt] at (9.8, 6.2) {SEL\_WB [1:0]};

    % 3. К DEST-DMUX, ACCA, ACCB, ALU: магистраль выходит из порта (11.4, 7.3)
    \draw[ctrlbus_line] (11.4, 7.3) -- (6.6, 7.3)
                        arc[start angle=0, end angle=180, radius=0.2]
                        -- (4.8, 7.3)
                        arc[start angle=0, end angle=180, radius=0.2]
                        -- (4.1, 7.3);
    
    % Вертикальный спуск к DEST-DMUX и ACCB по линии X = 4.1
    \draw[ctrlbus_line] (4.1, 7.3) -- (4.1, 4.9);
    \filldraw[red!80!black] (4.1, 7.3) circle (1.4pt);

    % Отвод к DEST-DMUX (SEL_ACC)
    \draw[ctrlbus] (4.1, 6.4) -- (3.35, 6.4);
    \filldraw[red!80!black] (4.1, 6.4) circle (1.4pt);
    \node[font=\tiny\bfseries, text=red!80!black, fill=white, inner sep=1pt, above] at (3.72, 6.45) {SEL\_ACC};

    % Отвод к ACCB (LD_B)
    \draw[ctrlbus] (4.1, 4.9) -- (ACCB.east);
    \node[font=\tiny\bfseries, text=red!80!black, fill=white, inner sep=1pt, above] at (3.9, 4.95) {LD\_B};

    % Магистраль продолжается влево над DEST-DMUX к ACCA и ALU:
    \draw[ctrlbus_line] (4.1, 7.3) -- (2.4, 7.3)
                        arc[start angle=0, end angle=180, radius=0.2]
                        -- (0.3, 7.3);

    % Отвод к ACCA (LD_A)
    \draw[ctrlbus] (0.3, 7.3) -- (0.3, 4.9) -- (ACCA.west);
    \node[font=\tiny\bfseries, text=red!80!black, fill=white, inner sep=1pt, left] at (0.25, 5.0) {LD\_A};

    % Отвод к ALU (ALU_OP [3:0])
    \draw[ctrlbus] (0.3, 4.9) -- (0.3, 2.8) -- (0.8, 2.8);
    \node[font=\tiny\bfseries, text=red!80!black, fill=white, inner sep=1pt, left] at (0.25, 3.2) {ALU\_OP [3:0]};

    % 4. К регистрам R0..R7: REG_WE [7:0]
    \draw[ctrlbus_line] (15.2, 9.2) -- (15.2, 12.8) -- (-0.2, 12.8);
    \node[font=\tiny\bfseries, text=red!80!black, fill=white, inner sep=1pt] at (13.0, 12.8) {REG\_WE [7:0]};
    \foreach \i in {0,...,7} {
        \draw[ctrlbus] (R\i_we |- 0, 12.8) -- (R\i_we);
        \filldraw[red!80!black] (R\i_we |- 0, 12.8) circle (1.4pt);
    }

\end{tikzpicture}
\end{document}